\chapter{Customer Needs}

\section{Session Overview}

Identifying customer needs is a key activity in the concept development phase of the product development process. The needs identified at this stage help guide the team in establishing product specifications, generating product concepts, and selecting a concept for further development.

Creating a high-quality information channel from customers to the product developers ensures that those who directly control the details of the product, including the product designers, fully understand the needs of the customer.



By the end of this session, you should be able to:
\begin{enumerate}
    \item generate effective survey or interview prompts;
    \item translate customer statements into clear needs statements;
    \item explain the value of lead users and extreme users; and
    \item compile a structured customer-needs table for your hydroponic product.
\end{enumerate}

\section{Key Terms}

The following terms are used throughout the chapter:
\begin{itemize}
    \item \textbf{Customer statement}: a direct quote or paraphrase describing a user's experience, frustration, or suggestion.
    \item \textbf{Needs statement}: a concise description of what the product must enable, reduce, or improve.
    \item \textbf{Lead user}: a user who experiences needs earlier than the mainstream market and may reveal emerging requirements.
    \item \textbf{Extreme user}: a user operating at the edge of normal use conditions and therefore likely to expose hidden limitations.
    \item \textbf{Variable type}: a broad category used to classify a need, such as functional, usability, cost, safety, or maintainability.
\end{itemize}

\section{Why Customer Needs Matter}

Customer needs help your group avoid designing only from assumptions. They provide a structured way to connect the mission statement from Sessions 5 and 6 to measurable engineering work. For hydroponic products, needs may involve fluid handling, maintenance, structural convenience, sensing accuracy, safety, or cost. The stronger your needs list, the easier it becomes to define meaningful specifications later.

\section{Using LLM for Customer Needs}

Given the quick timeline of our group assignment, it's essential to manage our expectations realistically. To fulfill the requirements for week 5's activity, you could, if feasible for your project, engage with potential customers to understand their needs. Alternatively, you may employ your creativity by envisioning yourself as the customer or utilize a Large Language Model (LLM) to generate possible customer needs. 


Should you opt for the LLM approach, please ensure to document/record your prompts and incorporate these prompts into your final submission, both in the ITEL and the whiteboard presentation. This documentation will be invaluable for both myself and your classmates.


\section{Pre-Class Activity 1: Generate Survey or Interview Prompts}

Before class, draft at least three questions or prompts that could help your group collect useful user feedback.

\subsection{Survey Design Guidelines}


To begin
\begin{enumerate}
    \item Consider the purpose and objectives of the survey. Such that, what specific information or insights you intend to obtain from the respondents. 
    For example, it can be about typical use, likes, dislike, or suggested improvement
    \item Based on the survey's goals, brainstorm potential questions and prompts that can help extract relevant and valuable responses from respondents.
    \item For this activity, select a atleast three questions that you believe would be most appropriate for the survey response.
    \begin{enumerate}
        \item Keep in mind that the question should align with the objectives and the value the question would provide. And it should be about typical use, likes, dislike, or suggested improvement
        \item It is beneficial if you can provide a justification explaining why you believe each of the question is appropriate for the survey response.
        \item Ensure that the questions are clear, concise, and unbiased to obtain accurate and meaningful survey responses.
    \end{enumerate}
    \item Note down the selected questions and their justifications.
\end{enumerate}
% Good questions should be:
% \begin{itemize}
%     \item clear and easy to understand;
%     \item neutral rather than leading;
%     \item connected to a real usage situation; and
%     \item useful for revealing frustrations, preferences, or desired improvements.
% \end{itemize}

Table~\ref{tab:session7-question-template} can be used to record your prompts.

\begin{table}[htbp]
\centering
\caption{Template for survey or interview prompts}
\label{tab:session7-question-template}
\renewcommand{\arraystretch}{1.2}
\begin{tabularx}{\textwidth}{|>{\raggedright\arraybackslash}p{0.08\textwidth}|>{\raggedright\arraybackslash}X|>{\raggedright\arraybackslash}p{0.28\textwidth}|}
\hline
\textbf{No.} & \textbf{Question or prompt} & \textbf{Why it is useful} \\
\hline
1 & How often do you check nutrient level or water level in your hydroponic system? & Reveals monitoring frequency and maintenance burden. \\
\hline
2 & What part of maintaining your hydroponic system is most frustrating? & Helps identify pain points and improvement opportunities. \\
\hline
3 & What would make the system easier to clean, refill, or monitor? & Encourages users to describe desired features directly. \\
\hline
\end{tabularx}
\end{table}

\section{Pre-Class Activity 2: Translate Customer Statements into Needs}

% After collecting observations, imagined use cases, or preliminary feedback, convert each customer statement into a needs statement.
\begin{enumerate}
    \item Imagine performing an everyday task using the chosen product but with an assumption that you encounter frustrations and difficulties performing the task. Based on this imaginative experience, identify the needs that the product's developer may have overlooked. Based on this experience, generate a customer statement.
    \item If time permits, observe someone else performing an everyday task using the chosen product and pay special attention to the frustrations and difficulties that person encounters.
    \begin{enumerate}
        \item Seek feedback from that person and generate a customer statement.
    \end{enumerate}
    \item Once you have generated the customer statements either by imaginative experience or feedback from a person, translate these customer statements into proper needs statements as per \ref{sec:Translation_guidance}.
\end{enumerate}

\subsection{Translation Guidance}
\label{sec:Translation_guidance}
A strong needs statement should:
    \begin{enumerate}
        \item Ensure that the needs statements are concise, specific, and clearly articulate the underlying requirement.
        \item Use clear and unambiguous language to accurately represent the customer's needs.
        \item Avoid any assumptions or biases when translating the customer statements into needs statements.
    \end{enumerate}

% Table~\ref{tab:session7-needs-example} shows a hydroponic example.

% \begin{table}[htbp]
% \centering
% \caption{Example of translating hydroponic customer statements into needs}
% \label{tab:session7-needs-example}
% \renewcommand{\arraystretch}{1.2}
% \begin{tabularx}{\textwidth}{|>{\raggedright\arraybackslash}X|>{\raggedright\arraybackslash}X|>{\raggedright\arraybackslash}p{0.2\textwidth}|}
% \hline
% \textbf{Customer statement} & \textbf{Needs statement} & \textbf{Variable type} \\
% \hline
% ``I never know when the nutrient tank is close to empty.'' & The system should provide a clear indication of nutrient level. & ??? \\
% \hline
% ``Cleaning the pipes takes too long.'' & The system should be easy to clean and maintain. & ??? \\
% \hline
% ``I worry that leaks will damage the floor near my plants.'' & The system should minimise leakage during normal operation. & ??? \\
% \hline
% ``I want something simple because I am new to hydroponics.'' & The system should be easy for a beginner to understand and operate. & ??? \\
% \hline
% \end{tabularx}
% \end{table}

\begin{table}[h!]
\centering
\begin{tabular}{|p{7cm}|p{7cm}|}
\hline
\textbf{Customer Statement} & \textbf{Need Statement} \\
\hline
“Why don’t you put protective shields around the battery contacts?” & The screwdriver battery is protected from accidental shorting \\
\hline
“I drop my screwdriver all the.” & The screwdriver operates normally after repeated dropping \\
\hline
“It doesn’t matter if it’s raining, I still need to work outside on Saturdays.” & The screwdriver operates normally in the rain \\
\hline
“I’d like to charge my battery from my Power backup.” & The screwdriver battery can be charged from a Power backup adapter \\
\hline
“I hate it when I don’t know how much juice is left in the batteries of cordless tools” & The screwdriver provides an indication of the energy level of the battery \\
\hline
\end{tabular}
\caption{Customer Statements and Corresponding Need Statements}
\end{table}


\section{Pre-Class Activity 3: Consider Lead Users and Extreme Users}

% Lead users and extreme users can reveal needs that ordinary users may not express clearly.

% \begin{itemize}
%     \item A \textbf{lead user} might be an experienced hydroponic grower who has already created workarounds for monitoring or dosing problems.
%     \item An \textbf{extreme user} might be someone operating in a very small apartment, in a very hot environment, or under strict cleaning constraints.
% \end{itemize}

% Write a short note explaining whether these user groups are important for your product and why. If they are relevant, identify what special insight they may provide.
The objective of this activity is reflect whether is essential to include lead users and extreme users in a survey related to your team's chosen product.

\smallskip
\noindent
To begin
\begin{enumerate}
    \item Reflect the objectives of the survey and the concept of lead users and extreme users.
    \item Consider the following points to determine whether it is essential to interview these two user groups based on objectives of the survey:
    \begin{enumerate}
        \item Lead Users:
        \begin{itemize}
            \item Assess whether lead users possess any attribute (e.g., unique insights, expertise, or experiences) that can contribute to understanding emerging needs and trends as well as the influence they may have on shaping future product developments or improvements.
        \end{itemize}
        \item Extreme Users:
        \begin{itemize}
            \item Assess whether extreme users can provide valuable insights into edge cases or challenging scenarios that may not occur in regular usage as well as their ability to help identify areas of improvement or innovation that cater to specific user segments.
        \end{itemize}
    \end{enumerate}
    \item Note down your analysis and justification.
\end{enumerate}

% \section{Optional Reflection: Needs and Innovation}

% Use this prompt for a short discussion or reflective paragraph: How can a carefully identified unmet need lead to a novel product concept? Which hidden or latent need in your hydroponic project might create the biggest design opportunity?

\section{In-Class Activity: Build the Customer Needs Table}

In class, merge your individual prompts and needs statements into one group table.

\subsection{Required Deliverable}

Your group should prepare a table containing:
\begin{itemize}
    \item the agreed survey or interview prompts;
    \item customer statements grouped by themes such as typical use, likes, dislikes, and suggested improvements;
    \item the interpreted needs statements; and
    \item the variable type for each interpreted need.
\end{itemize}

Table~\ref{tab:session7-group-template} provides the recommended format.



\begin{table}[htbp]
\centering
\caption{Compiled customer statements and interpreted needs for the selected hydroponic product}
\label{tab:session7-group-template}
\renewcommand{\arraystretch}{1.2}
\begin{tabularx}{\textwidth}{|>{\raggedright\arraybackslash}p{0.16\textwidth}|>{\raggedright\arraybackslash}X|>{\raggedright\arraybackslash}X|>{\raggedright\arraybackslash}p{0.16\textwidth}|}
\hline
\textbf{Prompt or theme} & \textbf{Customer statement} & \textbf{Interpreted need} & \textbf{Variable type} \\
\hline
Typical use &  &  &  \\
\hline
Like &  &  &  \\
\hline
Dislike &  &  &  \\
\hline
Suggested improvement &  &  &  \\
\hline
\end{tabularx}
\end{table}

Lecturer only


\begin{table}[h!]
\footnotesize
\centering
\caption{Customer Statements and Interpreted Needs for Thermostat Design}
\label{tab:interpreted_needs}
\begin{tabular}{|p{2.5cm}|p{5cm}|p{5cm}|p{2.5cm}|}
\hline
\textbf{Question/Prompt} & \textbf{Customer Statement} & \textbf{Interpreted Need} & \textbf{Variable Type} \\
\hline
Typical uses & I have to manually turn it on and off when it gets too hot or cold. & The thermostat maintains a comfortable temperature without requiring user action. & Functional \\
\cline{2-4}
& Each time I want to change the temperature, I need to adjust both thermostats in the house. & Any user input need not be made in multiple locations. (I) & Functional \\
\hline
Likes—current model & I like that I can change the temperature if the setting is too high. & The temperature setting is easy to control manually. & Usability \\
\cline{2-4}
& It didn’t cost a fortune. & The thermostat is affordable to purchase. & Cost \\
\hline
Dislikes—current model & I’m too lazy to learn how to figure it out. & The thermostat requires little or no user instruction or learning. & Usability \\
\cline{2-4}
& I sometimes forget to turn it off when we leave the house. & The thermostat saves energy when no one is home. (†) & Energy Saving \\
\cline{2-4}
& Sometimes the buttons don’t register a push and I have to press it repeatedly. & The thermostat definitively registers any user inputs. & Reliability \\
\hline
Suggested improvements & I would like my iPhone to adjust my thermostat. & The thermostat can be controlled remotely without requiring a special device. & Connectivity \\
\cline{2-4}
& I like to shift quickly between different options like energy saving or ultra comfort. & The thermostat responds immediately to differences in occupant preferences. & Responsiveness \\
\hline
\end{tabular}
\end{table}

While optional, you may create a new column at the end of the table, to indicate what type of need the interpreted needs (direct, latent, constant, variable, general, niche) is being addressed.


If your group uses an LLM to help organise responses, document the prompts used and verify that the grouped output still reflects the actual evidence collected.

\section{Summary and Next Steps}

At the end of this session, your group should have a clear customer-needs table and a short note about whether lead or extreme users matter for the project. In Session 8, these needs will be translated into measurable product specifications.
